\documentclass[11pt]{article}
% =====================================================================
%  Status:      research note (exploratory).  NOT a paper claim.
%  Provenance:  derived in the session of 2026-09-28/29 while reading
%               Paper 0 §§3--11 at commit d89ceb8; the source working
%               notes live outside this repository.
%  Promotion:   nothing here is promoted.  Items that could become paper
%               statements are listed in §6 and require an integration
%               audit under governance/05-mathematical-status.md §4.
% =====================================================================
\usepackage[utf8]{inputenc}
\usepackage[T1]{fontenc}
\usepackage{amsmath}
\usepackage{amssymb}
\usepackage{amsthm}
\usepackage{booktabs}
\usepackage{array}
\usepackage{hyperref}
\hypersetup{colorlinks=true,linkcolor=blue,urlcolor=cyan}
\usepackage[a4paper,margin=1in]{geometry}

\newtheorem{theorem}{Theorem}[section]
\newtheorem{lemma}[theorem]{Lemma}
\newtheorem{proposition}[theorem]{Proposition}
\newtheorem{corollary}[theorem]{Corollary}
\newtheorem{definition}[theorem]{Definition}
\newtheorem{remark}[theorem]{Remark}
\newtheorem{example}[theorem]{Example}

% Long file paths and inline formulas in \texttt cannot hyphenate; allow a
% little emergency stretch rather than printing overfull boxes.
\tolerance=1000
\emergencystretch=2.5em

\title{The joint quotient of the ACS charge and the AES affine operator}
\author{Mingli Yuan and DeepSeek Harness Agent}
\date{2026-09-29}

\begin{document}
\maketitle

\begin{abstract}
Paper 0 separates two readouts of a marked history: the charge endpoint, which
survives the abelianization of the history group, and the affine operator, which
is what the history actually does to the real line. This note studies the
\emph{coarsest quotient that keeps both}. That quotient is an explicit
three-dimensional solvable Lie group
$G=\mathbb R_A\times\operatorname{Aff}^+(1,\mathbb R)$, and it carries a natural
contact form whose non-integrability is exactly the failure of two arithmetic
steps to commute. The two readouts reappear as two projections of $G$, each
killing a different one-dimensional subgroup, and the pair recovers the group
element by a fibre-product condition. Everything in \S\S\ref{sec:group}--%
\ref{sec:contact} is proved here from the definitions; \S\ref{sec:boundaries}
records what is deliberately \emph{not} claimed.
\end{abstract}

\tableofcontents

\section{Status and provenance}
\label{sec:status}

\textbf{Status.} Research note. Exploratory. The propositions below are proved,
but the object $G$ is \emph{not} a Paper 0 object and nothing here is a paper
claim. The note is a candidate source for one; promotion requires the
integration audit of \texttt{governance/05-mathematical-status.md}~\S4.

\textbf{Provenance.} Written on 2026-09-29 from a reading of
\texttt{paper-0/} at commit \texttt{d89ceb8}. The reading notes and the exact
regression scripts behind the finite checks quoted here live outside this
repository. Where this note touches a claim that Paper 0 registers, the register
is authoritative and the conflict, if any, is recorded rather than reconciled
silently.

\textbf{Sources read for this note.} \path{paper-0/sections/03-aes-and-motion.tex}
(\texttt{def:regular-aes}, the canonical arithmetic frame, the four motions) and
\path{paper-0/sections/09-contact-of-tearing.tex} (the contact form, the
horizontal fields, the scalar horizontal differential).

\section{Two readouts of one history}
\label{sec:readouts}

Fix the two one-parameter families of arithmetic steps: the addition
$T_p:z\mapsto z+p$ and the positive dilation $D_s:z\mapsto sz$ with $s>0$. Write a
finite \emph{marked history} as a time-ordered word
$\gamma=(h_1,\dots,h_n)$ in these letters, and let
\[
\mathcal H=(\mathbb R,+)*(\mathbb R,+)
\]
be the free product generated by them, with the dilation generator parametrised by
$q=\log s$.

Paper 0 reads a history in (at least) two ways, and it is essential that they are
different readings of the same word rather than two objects in a chain:
\[
c(\gamma)=(A_\gamma,M_\gamma),
\qquad
\rho(\gamma)(z)=S_\gamma z+B_\gamma,\qquad S_\gamma=e^{M_\gamma}.
\]
Here $c$ is the \emph{charge endpoint}, obtained from the abelianization of
$\mathcal H$: $A_\gamma$ is the sum of the additive steps and $M_\gamma$ the sum
of the logarithmic dilations. And $\rho$ is the \emph{affine operator} actually
composed by the word.

The two are not functions of one another. The standard witness is
\begin{center}
\small
\begin{tabular}{@{}llccc@{}}
\toprule
history (time order) & note & $A$ & $S$ & $B$\\
\midrule
$T_1,\ D_2$                    &              & $1$ & $2$ & $2$\\
$D_2,\ T_1$                    &              & $1$ & $2$ & $1$\\
$D_2,\ T_2,\ D_{1/2},\ T_{-1}$ & identity op. & $1$ & $1$ & $0$\\
\bottomrule
\end{tabular}
\end{center}
The first two rows have equal charge and different operators; the third has the
identity operator and nonzero charge.

\begin{remark}
$c$ is the abelianization of the \emph{history} group. It is not the
abelianization of the affine group $\operatorname{Aff}^+(1,\mathbb R)$, which
would forget the translation part as well.
\end{remark}

\section{The joint quotient}
\label{sec:group}

\begin{definition}[Joint readout]
For a marked history $\gamma$, put
\[
J(\gamma)=\bigl(A_\gamma,\ S_\gamma,\ B_\gamma\bigr),\qquad S_\gamma>0 .
\]
\end{definition}

The joint readout is a homomorphism, and its image is an explicit group.

\begin{proposition}[Joint group]
\label{prop:joint-group}
Let
\[
G=\bigl\{(A,S,B):\ A,B\in\mathbb R,\ S>0\bigr\}
\]
with
\begin{equation}
\label{eq:law}
(A_2,S_2,B_2)(A_1,S_1,B_1)=(A_2+A_1,\ S_2S_1,\ B_2+S_2B_1),
\end{equation}
the right-hand side meaning \emph{first} $g_1$, \emph{then} $g_2$. Then
\begin{enumerate}
\item \eqref{eq:law} is a group law with unit $e=(0,1,0)$ and inverse
\[
(A,S,B)^{-1}=(-A,\ S^{-1},\ -B/S);
\]
\item $J$ extends to a surjective homomorphism $\mathcal H\twoheadrightarrow G$,
so $G$ is the image of $J$;
\item $G\cong\mathbb R\times\operatorname{Aff}^+(1,\mathbb R)$;
\item $G$ is the coarsest quotient of $\mathcal H$ through which both $c$ and
$\rho$ factor.
\end{enumerate}
\end{proposition}

\begin{proof}
(1) The third coordinate of $(g_3g_2)g_1$ and of $g_3(g_2g_1)$ is
$B_3+S_3B_2+S_3S_2B_1$ in both cases; the first two coordinates are immediate.
The stated inverse is checked by direct substitution.

(2) On generators, $J(T_p)=(p,1,p)$ and $J(D_s)=(0,s,0)$, and $J$ respects the
free-product product because $\rho$ is the composition and $c$ is additive over
concatenation. For surjectivity, let $s\ne1$ be any fixed dilation. For
$S\ne1$, choose $q=(B-A)/(S-1)$ and $p=A-q$; then
\[
(p,1,p)\,(0,S,0)\,(q,1,q)=(p+q,\ S,\ p+Sq)=(A,S,B).
\]
For $S=1$, choose $q=(B-A)/(s-1)$, $p=A-q$; then
\[
(p,1,p)\,(0,s,0)\,(q,1,q)\,(0,s^{-1},0)=(p+q,\ 1,\ p+sq)=(A,1,B).
\]
Each factor is the image of a generator or of its inverse, so every element of
$G$ is $J$ of some word.

(3) The map $(A,S,B)\mapsto(A,(S,B))$ is an isomorphism onto
$\mathbb R\times\operatorname{Aff}^+(1,\mathbb R)$, where
$\operatorname{Aff}^+(1,\mathbb R)$ is the group of maps $z\mapsto Sz+B$,
$S>0$, under composition.

(4) $J$ is surjective by (2), and any map out of $\mathcal H$ that factors both
$c$ and $\rho$ factors through $J$ by definition of $J$ as their pairing.
\end{proof}

\begin{remark}[$G$ is not the history]
\label{rem:not-history}
$J$ has a large kernel. The histories
\[
T_1,\ D_2,\ T_{-3},\ D_2,\ T_2
\qquad\text{and}\qquad
D_4
\]
are different reduced words in $\mathcal H$ --- five steps against one --- yet
both have $J=(0,4,0)$. No further projection out of $G$ can recover the step
count, the word, or the provenance.
\end{remark}

\begin{remark}[One useful family of realisers]
The history
\[
D_2,\ T_{2r},\ D_{1/2},\ T_{-r}
\]
realises the identity operator together with additive charge $r$:
$J=(r,1,0)$. Hence, having realised any affine operator, one may append such a
history to adjust $A$ independently of $(S,B)$.
\end{remark}

\subsection{Two characteristic subgroups}
\label{subsec:characteristic}

\begin{proposition}
\label{prop:center-derived}
The centre and the derived subgroup of $G$ are
\[
Z(G)=\{(A,0,0)\}\cong\mathbb R,
\qquad
[G,G]=\{(0,0,B)\}\cong\mathbb R ,
\]
and consequently
\[
G/[G,G]\cong\mathbb R^2,\qquad G/Z(G)\cong\operatorname{Aff}^+(1,\mathbb R).
\]
Both $Z(G)$ and $[G,G]$ are characteristic subgroups of $G$.
\end{proposition}

\begin{proof}
For a pure dilation $H_m=(0,m,0)$ and a pure offset $N_b=(0,0,b)$,
\[
H_mN_bH_m^{-1}N_b^{-1}=(0,0,(e^m-1)b),
\]
whose set of values is all of $\{(0,0,B)\}$; so the derived subgroup is exactly
$\{(0,0,B)\}$. An element commuting with every $N_b$ and every $H_m$ must have
$S=1$ (from the $H_m$) and $B=0$ (from the $N_b$), so the centre is exactly
$\{(A,0,0)\}$. Since the centre and the derived subgroup are preserved by every
automorphism, both are characteristic.
\end{proof}

\section{The two projections, and the fibre product}
\label{sec:projections}

Each readout is a projection of $G$.

\begin{proposition}[Charge projection]
\label{prop:pi-E}
Put $\pi_E(A,M,B)=(A,M)$. Then
\[
\pi_E(hg)=\pi_E(h)+\pi_E(g)\quad\text{additively in }\mathbb R^2,
\]
the action on $\mathbb R^2$ is by translations, the Euclidean metric
$g_E=du^2+dv^2$ is preserved, and
\[
\ker\pi_E=\{(0,0,B)\}=[G,G].
\]
\end{proposition}

\begin{proposition}[Operator projection]
\label{prop:pi-H}
Put $\pi_H(A,M,B)=(B,e^M)=(x,y)$, so that $y>0$. Then
\[
\pi_H(hg)=\bigl(b+e^m x,\ e^m y\bigr),
\]
the action is by positive affine maps of the upper half-plane, the hyperbolic
metric $g_H=(dx^2+dy^2)/y^2$ is preserved, the action is transitive, and
\[
\ker\pi_H=\{(A,0,0)\}=Z(G).
\]
\end{proposition}

\begin{proof}[Proof of both]
The formulas follow from \eqref{eq:law}. For $\pi_H$ metric preservation,
multiply numerator and denominator by $e^{-2m}$. Transitivity: given $x\in\mathbb
R$, $y>0$, take $B=x$, $M=\log y$. The kernels are read off from
\eqref{eq:law}.
\end{proof}

\begin{remark}[Two consistent conventions]
The earlier AES projection of Paper 0 is
$Q(A,M,B)=(-Be^{-M},\,e^{-M})$, which satisfies
\[
Q(g)=\pi_H(g^{-1}).
\]
The two are the left- and right-action conventions of one and the same
projection; neither replaces the other, and the assignment $a=-x/y$ of Paper 0 is
\emph{not} recovered from $\pi_H$.
\end{remark}

\begin{theorem}[Fibre product]
\label{thm:fibre-product}
$\pi_E$ and $\pi_H$ together determine the group element:
\[
G\ \cong\ \bigl\{((u,v),(x,y))\in\mathbb E^2\times\mathbb H^2:\ y=e^v\bigr\}
\ =\ \mathbb E^2\times_{\mathbb R}\mathbb H^2,
\]
the two maps to $\mathbb R$ being $v$ and $\log y$. The recovery formulas are
\[
A=u,\qquad M=v,\qquad B=x .
\]
\end{theorem}

\begin{proof}
Immediate: $\pi_H$ yields $x=B$ and $y=e^M$, and $v=M$ is the second coordinate
of $\pi_E$.
\end{proof}

Each projection forgets one direction, and \emph{which} direction it forgets is
not an artefact of the chosen coordinates.

\begin{proposition}
\label{prop:obstruction}
Any smooth homomorphism $G\to\operatorname{Isom}(\mathbb E^2)$ kills
$[G,G]=\{(0,0,B)\}$, and any transitive one has two-dimensional image, so its
kernel is exactly $[G,G]$. Likewise, any smooth homomorphism
$G\to\operatorname{Isom}(\mathbb H^2)$ with transitive image kills $Z(G)$.
\end{proposition}

\begin{proof}[Sketch]
Write $\mathfrak g=\operatorname{span}\{Z,H,N\}$ with $Z$ the $A$-direction, $H$
the $M$-direction, $N$ the $B$-direction, so that
\[
[Z,H]=[Z,N]=0,\qquad [H,N]=N .
\]
For the Euclidean case let $\phi(H)=(\omega,u)$ and $\phi(N)=(\eta,v)$ in
$\mathfrak e(2)=\mathfrak{so}(2)\oplus\mathbb R^2$, and let $J$ denote the
quarter-turn. Comparing the $N$ component of $[H,N]=N$ gives
$(\eta Jv-\omega Ju)=(0,v)$, hence $\eta=0$ and $\omega Jv=v$; pairing the latter
with $v$ gives $\|v\|^2=0$, so $\phi(N)=0$. For the hyperbolic case, a transitive
image is at least two-dimensional; the source is solvable and cannot map onto
$\mathfrak{sl}_2(\mathbb R)$, so the image is exactly two-dimensional and
non-abelian; hence $\phi(N)\ne0$, while $\phi(Z)$ commutes with both $\phi(H)$
and $\phi(N)$ and must therefore vanish.
\end{proof}

A companion metric statement, used in \S\ref{sec:contact}:

\begin{proposition}
\label{prop:left-invariant-metric}
The metric
\[
g_G=dA^2+dM^2+e^{-2M}dB^2
\]
is left-invariant on $G$, and substituting $x=B$, $y=e^M$ gives
\[
g_G=dA^2+\frac{dx^2+dy^2}{y^2},
\]
so $(G,g_G)$ is globally isometric to $\mathbb R\times\mathbb H^2$. Both
$\pi_E$ and $\pi_H$ are Riemannian submersions for suitable choices of vertical
directions, but their orthogonal horizontal spaces are different, and this metric
is a choice, not a metric determined by the group.
\end{proposition}

\begin{proof}
Left-invariance is the direct computation
$dA^2+dM^2+(s\,dB)^2/(sS)^2=dA^2+dM^2+dB^2/S^2$. The isometry statement is the
substitution. The final sentence is a warning rather than a claim: \S\ref{sec:contact}
uses a \emph{different} metric on the same group.
\end{proof}

\section{The contact form and the order defect}
\label{sec:contact}

Let $M=\log S$ and put
\begin{equation}
\label{eq:alpha}
\alpha=dB-dA-B\,dM .
\end{equation}

\begin{proposition}
\label{prop:contact}
With $X=\partial_A+\partial_B$ and $Y=\partial_M+B\partial_B$,
\[
\alpha(X)=\alpha(Y)=0,\qquad [X,Y]=\partial_B,\qquad
\alpha\wedge d\alpha=-dA\wedge dM\wedge dB\ne0 .
\]
Consequently the horizontal distribution $\mathcal D=\ker\alpha$ is maximally
non-integrable: the failure of two arithmetic steps to commute is the
non-integrability of $\mathcal D$.
\end{proposition}

\begin{proof}
Direct: $dB(X)=1$, $dA(X)=1$, $dM(X)=0$, so $\alpha(X)=0$; and $dB(Y)=B$,
$dA(Y)=0$, $dM(Y)=1$, so $\alpha(Y)=B-B=0$. The bracket is
$[\partial_A+\partial_B,\ \partial_M+B\partial_B]=\partial_B$. Finally
$d\alpha=dM\wedge dB$ and $B\,dM\wedge dM=0$, so
$\alpha\wedge d\alpha=-dA\wedge dM\wedge dB$.
\end{proof}

\begin{remark}
This is Paper 0's contact model in the multiplicative coordinate $S=e^M$; the
coordinates correspond to $u=A$, $v=M$, $a=B$. The note records it here only to
fix the coordinate conventions used by the companion notes.
\end{remark}

\begin{proposition}[Right translations preserve the data]
\label{prop:right-translation}
For fixed $g_0=(A_0,M_0,B_0)$, the right translation
\[
R_{g_0}(A,M,B)=(A+A_0,\ M+M_0,\ B+e^MB_0)
\]
satisfies $R_{g_0}^*\alpha=\alpha$ and preserves the horizontal metric
$dA^2+dM^2$. It induces translations on the charge projection and the hyperbolic
isometry $(x,y)\mapsto(e^{-M_0}(x-B_0),e^{-M_0}y)$ on the operator projection.
\end{proposition}

\begin{proof}
Substituting into \eqref{eq:alpha} gives
$d(B+e^MB_0)-d(A+A_0)-(B+e^MB_0)d(M+M_0)=\alpha$. The horizontal metric is
unchanged because $A,M$ are translated. The induced maps are the substitutions.
\end{proof}

\begin{remark}[Executing a step is not a symmetry]
\label{rem:not-symmetry}
Appending an arithmetic step is a \emph{left} multiplication in the convention
\eqref{eq:law}, and its flow need not preserve $\alpha$. Statement
\ref{prop:right-translation} is about structure-preserving right translations;
it does not say that the group generated by the arithmetic motions is the full
symmetry group of AES. The two questions must not be identified.
\end{remark}

\section{Boundaries: what is not claimed}
\label{sec:boundaries}

\begin{enumerate}
\item \textbf{$G$ is not a Paper 0 object.} Paper 0 owns the regular AES, the
canonical arithmetic frame, the four motions, the affine cocycle, the charge
endpoint, the contact structure and tearing. The joint quotient is a new
organising object introduced here; nothing in this note should be read as
attributing it to Paper 0.
\item \textbf{No finite-dimensionality claim for all of AES.} $G$ covers the
addition\,--\,positive-multiplication fragment only. Paper 0 already records that
the generators $A=\partial_a$, $M=a\partial_a$, $P=a\log a\,\partial_a$ do not
close into a finite-dimensional Lie algebra; the three-dimensional $G$ must not
be upgraded to a group carrying all arithmetic evolution.
\item \textbf{Discrete and syntactic levels are separate.} Discrete subgroups,
finite-word syntax, mixed-chirality marked histories and the original
computation histories are \emph{not} identified with this Lie group anywhere
above.
\item \textbf{No logarithms in finite coordinates.} In the multiplicative
coordinate $(A,S,B)$ the group law \eqref{eq:law} uses no logarithm. Statements
that pair an $M$-coordinate predicate (for instance $A=M$) with the same object
must keep the relation $M=\log S$ explicitly; it does not become polynomial by
change of coordinates.
\item \textbf{Exactness of the finite checks.} Every numerical statement quoted
in the companion notes is either an identity proved above or an exact rational
regression; no floating-point value enters any judgement.
\end{enumerate}

\begin{thebibliography}{9}
\bibitem{Paper0}
M.~Yuan and DeepSeek Harness Agent,
\emph{Arithmetic Expression Geometry 0: Arithmetic Expression Spaces},
\texttt{paper-0/}, commit \texttt{d89ceb8f8090482317b750407f623c8dcfb37ea2}.
\end{thebibliography}

\end{document}
